\begin{longtable}{|p{1.5cm}|p{5cm}|p{7cm}|}
\caption{Kebutuhan Fungsional}
\label{tbl:kebutuhan-fungsional} \\
\hline
\textbf{ID} & \textbf{Kebutuhan} & \textbf{Penjelasan} \\
\hline
\endfirsthead

\hline
\textbf{ID} & \textbf{Kebutuhan} & \textbf{Penjelasan} \\
\hline
\endhead

\hline
\endfoot

FR-01 &
Sistem dapat melakukan proses \textit{crawling} data secara \textit{stealth} &
Sistem harus memfasilitasi metode \textit{crawling} bertahap (\textit{low and slow}), 
memanfaatkan proses atau API sah (\textit{living off the land}) agar pola akses menyerupai aktivitas normal 
sehingga sulit terdeteksi \textit{anti-malware}. \\
\hline

FR-02 &
Sistem mendukung eksfiltrasi data menggunakan metode yang sulit dianalisis &
Eksfiltrasi dilakukan melalui kanal terenkripsi, 
data dicicil secara acak, serta menyamarkan trafik sehingga tidak mudah dibedakan dari aktivitas jaringan normal oleh deteksi keamanan. \\
\hline

FR-03 &
Sistem tidak menyimpan artefak proses \textit{crawling} atau eksfiltrasi di disk &
Semua proses berjalan \textit{in-memory}, tanpa membuat file \textit{temporary} yang dapat dianalisis secara forensik, 
serta mendukung \textit{memory wiping} setelah selesai eksekusi. \\
\hline

FR-04 &
Sistem memungkinkan penyembunyian aktivitas \textit{crawling} dan eksfiltrasi dalam proses aplikasi sah &
Aktivitas bisa disisipkan melalui teknik \textit{injection/hollowing} ke aplikasi yang dipercaya, 
agar \textit{event crawling} dan eksfiltrasi tersamarkan sebagai aktivitas aplikasi legal. \\
\hline

FR-05 &
Sistem dapat melakukan \textit{clean up} otomatis untuk menghilangkan jejak digital &
Setelah aktivitas \textit{crawling} dan eksfiltrasi selesai, sistem menghapus artefak, jejak \textit{log}, atau \textit{payload} sementara 
sehingga \textit{anti-malware} maupun forensik tidak menemukan indikasi aktivitas berbahaya. \\
\hline

\end{longtable}